\documentclass[12pt,a4paper]{report}

\usepackage[T1]{fontenc}
\usepackage[utf8]{inputenc}
\usepackage[english]{babel}
\usepackage{lmodern}
\usepackage{geometry}
\usepackage{setspace}
\usepackage{hyperref}
\usepackage{enumitem}

\geometry{margin=2.5cm}
\onehalfspacing
\setlist[itemize]{noitemsep, topsep=4pt}

\hypersetup{
  colorlinks=true,
  linkcolor=blue,
  urlcolor=blue,
  pdftitle={PFE Internship Report - Chapter 1},
  pdfauthor={Your Name}
}

\begin{document}

% =====================
% Optional simple title page
% =====================
\begin{titlepage}
  \centering
  {\Large Internship Report (PFE)\par}
  \vspace{1.5cm}
  {\huge \textbf{Collectra Project}\par}
  \vspace{0.5cm}
  {\Large Chapter 1: General Presentation and Project Framing\par}
  \vfill

  \begin{flushleft}
    \textbf{Student:} Your Full Name \\
    \textbf{Institution:} Your Institution Name \\
    \textbf{Host Organization:} Company / Lab Name \\
    \textbf{Academic Year:} 2025--2026
  \end{flushleft}

  \vfill
  {\large \today\par}
\end{titlepage}

\tableofcontents
\clearpage

\chapter{General Presentation and Project Framing}
\label{chap:general-presentation}

\section*{Introduction}
Digital transformation in commercial activities requires companies to structure how they track customers, debts, and collection actions. In this context, the quality of internal tools has a direct impact on team responsiveness, operational traceability, and data reliability for decision-making.

This chapter introduces the overall framework of the \textbf{Collectra} project, developed during a Final Year Internship (PFE). It presents the context, the core problem, project objectives, functional scope, and the high-level technical choices that guide the rest of the report.

\section{General Context}
Debt collection and follow-up operations are often handled through heterogeneous processes (spreadsheets, manual exchanges, and disconnected tools), which typically leads to:
\begin{itemize}
  \item scattered customer information;
  \item weak visibility over action history;
  \item coordination errors between team members;
  \item limited real-time insight into case progress.
\end{itemize}

The \textbf{Collectra} project addresses this need through a SaaS web platform built around multi-workspace organization, with secure access and standardized follow-up operations.

\section{Collectra Project Overview}
\subsection{Project Nature}
Collectra is a full-stack web application designed to centralize the management of:
\begin{itemize}
  \item customers;
  \item debts;
  \item payment promises;
  \item action timelines related to collection cases;
  \item team members and roles within each workspace.
\end{itemize}

The application follows a modern architecture with clear separation between frontend, backend API, and data access layers.

\subsection{Functional Positioning}
The platform is designed to provide:
\begin{itemize}
  \item a private dashboard for authenticated users;
  \item workspace management to isolate data by tenant context;
  \item team administration (invitation, role assignment, activation/deactivation);
  \item secure and documented APIs for business operations.
\end{itemize}

\section{Problem Statement}
The central problem can be formulated as follows:

\begin{quote}
\textit{How can we design a secure, multi-tenant, and scalable web platform that efficiently supports the debt follow-up lifecycle while ensuring action traceability and strict data isolation across workspaces?}
\end{quote}

This problem includes both technical challenges (security, architecture, scalability) and organizational challenges (collaboration, role governance, continuity of follow-up).

\section{Project Objectives}
\subsection{General Objective}
Develop a SaaS solution named \textbf{Collectra} to centralize customer and debt management with robust access control and multi-workspace organization.

\subsection{Specific Objectives}
The main operational objectives are:
\begin{itemize}
  \item implement secure authentication (sessions, tokens, protected APIs);
  \item implement workspace management and tenant-context resolution;
  \item provide business modules for customers, debts, promises, and actions;
  \item support team management (invitations, \texttt{MANAGER}/\texttt{AGENT} roles, member status);
  \item ensure data consistency through a relational database and type-safe ORM;
  \item provide a maintainable, testable, and extensible software foundation.
\end{itemize}

\section{Scope of the Project}
\subsection{In-Scope Features}
The current functional scope includes:
\begin{itemize}
  \item authentication and user account management;
  \item workspace management and active workspace selection;
  \item business entities (customers, debts, promises, actions);
  \item internal member and invitation management;
  \item a web interface for operational use.
\end{itemize}

\subsection{Out-of-Scope at This Stage}
At this stage, the following items are not part of the core deliverable:
\begin{itemize}
  \item advanced integrations with external ERP systems;
  \item full business intelligence analytics;
  \item multi-region deployment at very large scale.
\end{itemize}

\section{High-Level Technology Overview}
The project is implemented in a monorepo architecture with the following stack:
\begin{itemize}
  \item \textbf{Frontend:} Next.js (App Router), React, TypeScript, Tailwind CSS;
  \item \textbf{Backend:} Hono.js, TypeScript, Zod validation, OpenAPI documentation;
  \item \textbf{Data Layer:} PostgreSQL (Supabase) and Prisma ORM;
  \item \textbf{Auth \& Security:} Supabase Auth, JWT, HTTP-only cookies, middleware-based access control;
  \item \textbf{Code Organization:} Turborepo with pnpm workspaces (apps + shared packages).
\end{itemize}

This foundation promotes clear responsibility boundaries, code reuse, and progressive feature evolution.

\section{Development Approach (Overview)}
The project follows an incremental implementation approach:
\begin{itemize}
  \item analyze priority functional requirements;
  \item model data and apply migration workflows;
  \item implement secure business APIs;
  \item integrate frontend modules and state management;
  \item validate features through targeted technical testing.
\end{itemize}

This method enables iterative delivery of usable functionality while maintaining code quality and system coherence.

\section*{Conclusion}
This first chapter introduced the overall framework of the Collectra project, including its context, problem statement, objectives, and scope. It also summarized the key technical choices and the development methodology.

The next chapter will present the requirement analysis and system design, focusing on actors, use cases, and data modeling.

\end{document}
